\section{Partie Théorique}
\subsection{Calculs d'inertie}
Afin de pouvoir déterminer les couples nécessaires à la réalisation des profils de vitesse imposés, il est indispensable d’évaluer l’inertie équivalente du système entraîné par le moteur pas-à-pas. Cette inertie totale rapportée au moteur, notée $J_{m,\text{tot}}$, résulte de la contribution des moteurs, du codeur, ainsi que de l’ensemble des éléments mécaniques en rotation dans la chaîne de transmission.

Les données relatives aux moteurs, au codeur incrémental et à la transmission sont extraites des fiches techniques fournies et résumées dans le Tableau~\ref{tab :donnees_moteur}. Les dimensions géométriques et les matériaux des différentes pièces mécaniques de l’axe principal sont quant à eux donnés dans le Tableau~\ref{tab :Geometrie_inertie}.
\begin{table}[h!]
\centering
\begin{tabular}{|l| c| c|}
\hline
Description & Unité & Valeur \\
\hline
Inertie moteur pas-à-pas & g·cm$^2$ & 77 \\
\hline
Inertie moteur DC & kg·m$^2$ & $6.1 \times 10^{-7}$ \\
\hline
Facteur de tension $k_u$ moteur DC & V / 1000 rpm & 7.2 \\
\hline
Inertie codeur moteur DC & g·cm$^2$ & 0.4 \\
\hline
Nombre de dents petite roue & -- & 15 \\
\hline
Nombre de dents grande roue & -- & 75 \\
\hline
Rapport de transmission $i$ & -- & 5 \\
\hline
Rendement $\eta$ & \% & 88.5 \\
\hline
\end{tabular}
\caption{Données des moteurs et de la transmission}
\label{tab :donnees_moteur}
\end{table}
\begin{table}[h!]
\centering
\begin{tabular}{|l| c| c| c| c| c|}
\hline
Élément & Matériau & $R_{ext}$ [mm] & $R_{int}$ [mm] & $H$ [mm] & $\rho$ [kg/m$^3$] \\
\hline
Plateau principal & PMMA & 95 & 3 & 5 & 1880 \\
 & Alu & 6 & 3 & 10 & 2700 \\
 & Alu & 20 & 3 & 5 & 2700 \\
 \hline
Axe vertical & Acier & 3 & -- & 55 & 7850 \\
 & Acier & 5 & -- & 60 & 7850 \\
 \hline
Roue dentée (2x) & Alu & 29.5 & 3 & 9 & 2700 \\
 & Alu & 10 & 3 & 5 & 2700 \\
 \hline
Roulements (2x) & Acier & 15 & 8 & 7 & 7850 \\
\hline
Plateau secondaire & PMMA & 20 & 3 & 2 & 1880 \\
 & Alu & 15 & 3 & 3 & 2700 \\
 & Alu & 7.5 & 3 & 8 & 2700 \\
\hline
\end{tabular}
\caption{Caractéristiques géométriques des éléments de l’axe principal}
\label{tab :Geometrie_inertie}
\end{table}
\subsubsection{Inerties des moteurs et du codeur}

Les inerties propres au moteur pas-à-pas, au moteur à courant continu utilisé comme capteur de vitesse, ainsi qu’au codeur incrémental sont extraites des fiches techniques fournies. L’inertie du moteur DC est convertie dans des unités cohérentes avec le reste des calculs. En tenant également compte de l’inertie de la petite roue dentée montée sur l’axe du moteur, les inerties des axes moteurs s’écrivent :

\begin{equation}
J_{DC} = J_{\text{mot,DC}} + J_{\text{codeur}} + J_{\text{roue}} = 6.92 \; \text{g·cm}^2
\end{equation}

\begin{equation}
J_{\text{PàP}} = J_{\text{mot,PàP}} + J_{\text{roue}} = 77.42 \; \text{g·cm}^2
\end{equation}

\subsubsection{Inerties des éléments mécaniques}

Les éléments mécaniques entraînés par l’axe principal comprennent les plateaux, l’axe vertical, les roues dentées et les roulements. Ces composants possèdent des masses non négligeables et doivent donc être pris en compte dans le calcul de l’inertie globale.

Chaque pièce est modélisée soit comme un cylindre plein, soit comme un cylindre creux selon sa géométrie. Les inerties sont calculées à l’aide des relations suivantes :

\begin{equation}
J_{\text{cylindre plein}} = \frac{1}{2} m R_{\text{ext}}^2
\end{equation}

\begin{equation}
J_{\text{cylindre creux}} = \frac{1}{2} m \left( R_{\text{ext}}^2 + R_{\text{int}}^2 \right)
\end{equation}

Les masses sont déterminées à partir des dimensions géométriques et des masses volumiques des matériaux indiquées dans le Tableau~2. La somme des inerties de l’ensemble des éléments de l’axe principal conduit à une inertie totale :

\begin{equation}
J_{\text{axe}} = 12745.7 \; \text{g·cm}^2
\end{equation}

\subsubsection{Réduction des inerties au moteur pas-à-pas}

Le système comporte plusieurs axes reliés par des transmissions à courroie. Afin de déterminer l’inertie équivalente vue par le moteur pas-à-pas, les inerties des différents axes sont ramenées à cet axe en tenant compte du rapport de transmission $i = 5$ et du rendement $\eta = 0.885$.

L’inertie totale ramenée au moteur pas-à-pas s’exprime alors comme :

\begin{equation}
J_{m,\text{tot}} =
J_{\text{PàP}}
+ \frac{J_{\text{axe}}}{\eta \, i^2}
+ \frac{J_{DC}}{\eta^2 \, i^4}
\end{equation}

En remplaçant par les valeurs numériques obtenues :

\begin{equation}
J_{m,\text{tot}} =
77.42
+ \frac{12745.7}{0.885 \cdot 25}
+ \frac{6.92}{0.885^2 \cdot 625}
\end{equation}

On obtient finalement :

\begin{equation}
\boxed{
J_{m,\text{tot}} \approx 662.3 \; \text{g·cm}^2
}
\end{equation}

Cette inertie équivalente sera utilisée dans la suite du laboratoire pour le calcul des couples moteurs nécessaires lors des phases d’accélération et de décélération.
\subsection{Profils de vitesse}
Dans cette section, les profils de vitesse utilisés lors de la partie expérimentale sont analysés et leurs caractéristiques théoriques sont déterminées. Trois profils de vitesse trapézoïdaux symétriques sont considérés. Chaque profil est défini par une phase d’accélération à accélération constante, une phase à vitesse maximale constante, puis une phase de décélération à accélération constante opposée.  

Pour tous les profils étudiés, la durée totale du mouvement ainsi que le déplacement angulaire total sont identiques. Les valeurs imposées sont les suivantes :
\begin{equation}
t_{\text{mouv}} = 0.75 \; \text{s}
\qquad
\Theta_{\text{plateau}} = 90^\circ
\end{equation}

En tenant compte du rapport de transmission entre le plateau principal et le moteur pas-à-pas, égal à $i = 5$, le déplacement angulaire total du moteur s’écrit :
\begin{equation}
\Theta_{\text{tot}} = i \cdot \Theta_{\text{plateau}}
= 5 \cdot 90^\circ
= 450^\circ
= \frac{5\pi}{2} \; \text{rad}
\end{equation}

Les profils de vitesse sont décrits à l’aide de fractions $p_1$, $p_2$ et $p_3$ correspondant respectivement aux phases d’accélération, de vitesse constante et de décélération, exprimées en fraction de la durée totale du mouvement. La répartition des profils étudiés est donnée dans le Tableau~\ref{tab:profils}.

\begin{table}[h]
\centering
\renewcommand{\arraystretch}{1.5}
\begin{tabular}{|c|c|c|c|}
\hline
 & $p_1$ & $p_2$ & $p_3$ \\
\hline
Profil 1 & $\frac{1}{3}$ & $\frac{1}{3}$ & $\frac{1}{3}$ \\
\hline
Profil 2 & $\frac{1}{4}$ & $\frac{1}{2}$ & $\frac{1}{4}$ \\
\hline
Profil 3 & $\frac{1}{8}$ & $\frac{3}{4}$ & $\frac{1}{8}$ \\
\hline
\end{tabular}
\caption{Répartition des profils de vitesse}
\label{tab:profils}
\end{table}


Ces profils étant symétriques, l’accélération maximale est égale en valeur absolue à la décélération maximale. Les durées des différentes phases s’écrivent :
\begin{equation}
t_1 = p_1 T, \qquad
t_2 = p_2 T, \qquad
t_3 = p_3 T
\end{equation}

où $T = t_{\text{mouv}}$.

Le déplacement angulaire total peut être obtenu en intégrant la vitesse sur toute la durée du mouvement. Pour un profil trapézoïdal symétrique, on obtient :
\begin{equation}
\Theta_{\text{tot}} =
\frac{1}{2}\Omega_{\max} t_1
+ \Omega_{\max} t_2
+ \frac{1}{2}\Omega_{\max} t_3
\end{equation}

ce qui peut se réécrire sous la forme :
\begin{equation}
\Theta_{\text{tot}} =
\Omega_{\max}
\left(
-\frac{1}{2}t_3 + t_3 + t_2 + \frac{1}{2}t_1
\right)
\end{equation}

La vitesse maximale du moteur est alors donnée par :
\begin{equation}
\Omega_{\max} =
\frac{\Theta_{\text{tot}}}
{
-\frac{1}{2}t_3 + t_3 + t_2 + \frac{1}{2}t_1
}
\end{equation}

L’accélération maximale est obtenue à partir de la relation de mouvement circulaire uniformément accéléré :
\begin{equation}
\alpha_{\max} = \frac{\Omega_{\max}}{t_1}
\end{equation}

La vitesse angulaire moyenne du moteur, identique pour les trois profils puisque le déplacement total et la durée sont constants, est définie par :
\begin{equation}
\Omega_{\text{moy}} = \frac{\Theta_{\text{tot}}}{T}
\end{equation}

En remplaçant par les valeurs numériques, on obtient :
\begin{equation}
\Omega_{\text{moy}} =
\frac{5\pi/2}{0.75}
= 10.47 \; \text{rad·s}^{-1}
\end{equation}

Les valeurs de vitesse maximale, d’accélération maximale et de vitesse moyenne obtenues pour chacun des profils sont regroupées dans le Tableau~\ref{tab:valeurs_profils}.

\begin{table}[h]
\centering
\begin{tabular}{|c|c|c|c|}
\hline
Profil & $\Omega_{\max}$ [rad/s] & $\alpha_{\max}$ [rad/s$^2$] & $\Omega_{\text{moy}}$ [rad/s] \\
\hline
Profil 1 (1/3--1/3--1/3) & 15.71 & 62.8 & 10.47 \\
Profil 2 (1/4--1/2--1/4) & 13.97 & 74.5 & 10.47 \\
Profil 3 (1/8--3/4--1/8) & 11.97 & 127.7 & 10.47 \\
\hline
\end{tabular}
\caption{Valeurs théoriques des profils de vitesse}
\label{tab:valeurs_profils}
\end{table}

Ces résultats montrent que, bien que la vitesse moyenne reste constante pour tous les profils, les accélérations maximales augmentent lorsque la durée des phases d’accélération et de décélération diminue. Ces accélérations auront une influence directe sur le couple maximal requis au niveau du moteur pas-à-pas, qui sera analysé dans la section suivante.



\subsection{ Calcul des valeurs de programmation}
Afin de piloter correctement le moteur pas-à-pas pour réaliser les profils de mouvement calculés précédemment, il est nécessaire de déterminer les paramètres à programmer dans l’électronique de commande. Ces paramètres incluent notamment : \texttt{velocity}, \texttt{amax}, \texttt{Position} (ou \texttt{Temps}), \texttt{pulse\_div}, \texttt{ramp\_div} et \texttt{Usrs}.

\subsubsection{Fréquence des micro-pas}

La première étape consiste à calculer la fréquence des micro-pas \(f_{\text{usf}}\) [Hz], qui dépend de la résolution des micro-pas définie par \texttt{Usrs}. Pour notre cas, la résolution choisie est :
\[
\text{Usrs} = 6 \quad \Rightarrow \quad 64 \text{ micro-pas par pas entier}.
\]

La fréquence des micro-pas est calculée à partir de la fréquence de l’horloge interne \(f_{\text{CLK}}\), du nombre de pas par tour \(N\) et de la vitesse de pas entière (\texttt{velocitymax}) selon la relation suivante :
\begin{equation}
f_{\text{usf}} = \frac{N \cdot \text{velocitymax} \cdot 2}{\text{Usrs}}
\end{equation}

\subsubsection{Pré-diviseur de vitesse \texttt{pulse\_div}}

Le paramètre \texttt{pulse\_div} permet d’adapter la vitesse maximale du moteur à la fréquence des micro-pas. Il est calculé par :
\begin{equation}
\text{pulse\_div} = \log_2 \left( \frac{f_{\text{CLK}} \cdot \text{velocitymax}}{f_{\text{usf}} \cdot 2048 \cdot 32} \right)
\end{equation}
Cette valeur est arrondie à l’entier inférieur. Elle est ensuite utilisée pour recalculer la vitesse maximale réelle à programmer :
\begin{equation}
\text{velocity} = \frac{f_{\text{usf}} \cdot 2^{\text{pulse\_div}}}{2048 \cdot 32} \cdot \frac{1}{f_{\text{CLK}}}
\end{equation}
La vitesse est arrondie à l’entier supérieur tant que sa valeur reste inférieure à \texttt{velocitymax}.

\subsubsection{Accélération et pré-diviseur de rampe \texttt{ramp\_div}}

L’accélération angulaire maximale en micro-pas est calculée à partir de l’accélération maximale du moteur (\(\alpha_{\max}\)) et de la résolution des micro-pas :
\begin{equation}
a = \text{résolution} \cdot \frac{\alpha_{\max}}{2\pi} \; [\text{Hz/s}]
\end{equation}

Le paramètre \texttt{ramp\_div} est ensuite calculé comme :
\begin{equation}
\text{ramp\_div} = \log_2 \left( \frac{f_{\text{CLK}}^2 \cdot \text{amax}}{a} \right) - \text{pulse\_div} - 29
\end{equation}

La valeur arrondie à l’entier inférieur permet ensuite de recalculer la valeur effective de l’accélération à programmer :
\begin{equation}
a = \frac{f_{\text{CLK}}^2 \cdot \text{amax}}{2^{\text{pulse\_div} + \text{ramp\_div} + 29}}
\end{equation}
Cette valeur est arrondie à l’entier supérieur tant qu’elle reste inférieure à \texttt{amax}.

\subsubsection{Consigne de position et temps}

Si le moteur est contrôlé en position angulaire, la consigne \(\vartheta_c\) [pas] est calculée par :
\begin{equation}
\vartheta_c = \frac{N \cdot 2}{\text{Usrs}} \cdot \frac{\vartheta}{360}
\end{equation}
où \(\vartheta\) est l’angle cible en degrés.  

Si le moteur est contrôlé en temps, le paramètre \texttt{Temps} correspond à la durée totale du mouvement exprimée en tranches de 10 ms :
\begin{equation}
\text{Temps} = T[\text{s}] \cdot 100
\end{equation}

\subsubsection{Variables à programmer}

Le Tableau~\ref{tab:variables_prog} résume les principales variables à entrer dans le programme de l’électronique pour obtenir le profil de mouvement souhaité.

\begin{table}[h]
\centering
\renewcommand{\arraystretch}{1.3}
\begin{tabular}{|c|c|c|c|c|c|}
\hline
Nom & Description & N° param. & Gamme & Valeur par défaut & Unité \\
\hline
fCLK & Fréquence de l’horloge & --- & --- & 16 & MHz \\
velocity & Vitesse maximale & 4 & 0 … 2047 & 1000 & [-] \\
amax & Accélération maximale & 5 & 0 … 2047 & 1000 & [-] \\
pulse\_div & Pré-diviseur vitesse & --- & 0 … 13 & 3 & [-] \\
ramp\_div & Pré-diviseur rampe & 153 & 0 … 13 & 7 & [-] \\
Usrs & Résolution micro-pas & 140 & 0 … 6 & 6 & [-] \\
N & Nombre de pas par tour & --- & --- & 200 & pas/tr \\
\hline
\end{tabular}
\caption{Variables de programmation pour l’électronique du moteur}
\label{tab:variables_prog}
\end{table}

Le Tableau~\ref{tab:Usrs} donne la correspondance entre la valeur de \texttt{Usrs} et le nombre réel de micro-pas par pas entier.

\begin{table}[h]
\centering
\renewcommand{\arraystretch}{1.3}
\begin{tabular}{|c|c|}
\hline
Usrs & Nombre de micro-pas \\
\hline
0 & 1 (pas entier) \\
1 & 2 (demi-pas, non recommandé) \\
2 & 4 \\
3 & 8 \\
4 & 16 \\
5 & 32 \\
6 & 64 \\
7 & 128 \\
8 & 256 \\
\hline
\end{tabular}
\caption{Nombre de micro-pas en fonction de la valeur de Usrs}
\label{tab:Usrs}
\end{table}

Ces calculs permettent de programmer correctement la vitesse, l’accélération et la position (ou le temps) afin que le moteur exécute les profils de mouvement souhaités de manière précise et fiable. Les ajustements éventuels de \texttt{pulse\_div} et \texttt{ramp\_div} peuvent être faits par incréments unitaires pour optimiser le comportement réel du moteur.

